\chapter{Wie eine vollständige Rekonstruktion aussehen könnte}
\label{ch:recon}
\begin{bild}{Fortschritt v1.3}
Ein gemeinsamer endlicher Träger realisiert nun Aufzeichnung, Näherungspräparation, Reset und Mehrzeitantwort unter einem ausdrücklich angegebenen Ressourcenvertrag. Das beantwortet eine bedingte Existenzfrage. Es wählt weder den Vertrag aus TFPT allein noch ein universelles positives Zustandsfunktional aus. Die nachfolgenden Anforderungen bleiben deshalb nötig; die Konstruktion steht in Kapitel~\ref{ch:execution13}.
\end{bild}
\section{Vom Wunsch nach Eindeutigkeit zu einer prüfbaren Frage}
Eine endliche Liste von Schatten bestimmt im Allgemeinen mehrere Prozesse. Man kann sogar jeden Kandidaten um einen ungekoppelten unsichtbaren Hilfsfaktor ergänzen, auf den alle erlaubten Experimente trivial wirken. Alle Antworten bleiben gleich. Eindeutigkeit ist daher zunächst nur für die minimal beobachtbare Darstellung erreichbar.

Die operationelle Gleichheit aus \eqref{eq:op} nimmt genau diesen Gedanken ernst. Die Klasse erlaubter Experimente muss unter Fortsetzung und zulässiger Einbettung in größere Experimente geschlossen sein. Erst dann kann man verschiedene Beschreibungen konsistent zusammenfassen. Das ist keine Behauptung, die Welt bestehe aus klassischen Geschichten; es ist eine Regel zur Bewertung ihrer unterscheidbaren Antworten.

Im endlichen linearen Fall ist sie berechenbar. Man startet mit dem Raum der Ausleseeffekte und erweitert ihn durch die rückwärts wirkenden erlaubten Operationen:
\begin{equation}
V_0=\operatorname{span}\{\text{Ausleseeffekte}\},\qquad
V_{k+1}=V_k+\operatorname{span}\{\mathcal I^*(E):E\in V_k\}.
\end{equation}
Wenn der Raum stabil wird, enthält er alle linearen Zukunftstests. Zustände, die auf ihm dieselben Werte liefern, sind relativ zu diesem Zugriff gleich. Die 30- und 45-dimensionalen Beispiele sind konkrete Fälle dieser Methode. \src{S05, §11; S06}

\section{Eine kleinere Beschreibung ist nicht automatisch eine Algebra}
Eine Projektion kann wichtige Wahrscheinlichkeiten erhalten und trotzdem die Multiplikation zerstören. S03 gibt einen einfachen Gegenbeleg: Projiziere $M_2(\C)$ auf $\operatorname{span}\{I,X,Z\}$ und definiere $A\star B=P(AB)$. Dann
\begin{equation}
(X\star X)\star Z=Z,\qquad X\star(X\star Z)=0.
\end{equation}
Die neue Multiplikation ist nicht assoziativ. Eine nützliche reduzierte Vorhersagebeschreibung darf also nicht ohne zusätzlichen Nachweis als geschlossene physische Operationsalgebra verwendet werden.

\section{Was vollständige Phaseninformation zusätzlich bringt}
Eine phasentreue Rekonstruktion braucht mehr als eine Tabelle einzelner Wahrscheinlichkeiten. Ein passender Ausgangspunkt ist eine komplexe Algebra $\mathcal A$ mit Einheit, Multiplikation, Adjunktion und einem positiven normierten Zustandsfunktional $\omega$:
\begin{equation}
\omega(1)=1,\qquad\omega(a^\dagger a)\ge0.
\end{equation}
Für Operationswörter $u,v$ definiert
\begin{equation}
G(u,v)=\omega(u^\dagger v)
\end{equation}
einen positiven semidefiniten Korrelationskern. Er enthält die relativen komplexen Phasen zwischen verschiedenen Vorgängen.

Die GNS-Konstruktion nimmt die Nullrichtungen
\begin{equation}
\mathcal N=\{a:\omega(a^\dagger a)=0\},\qquad
\mathcal H_\omega=\overline{\mathcal A/\mathcal N}
\end{equation}
heraus und stellt Operationen durch Linksanwendung dar: $\pi(a)[b]=[ab]$. Der zyklische Zustand ist $\Omega_\omega=[1]$. Dann gilt
\begin{equation}
\omega(a)=\langle\Omega_\omega,\pi(a)\Omega_\omega\rangle.
\end{equation}
Für eine vollständig festgelegte C*-Algebra mit Zustand ist die minimale zyklische Darstellung bis auf unitäre Äquivalenz bestimmt. Bei unbeschränkten Generatoren kommen Domänen- und Stetigkeitsfragen hinzu.

\begin{bild}{Eine vollständige Partitur rekonstruiert die Aufführungsklasse}
Kennt man nicht nur die Lautstärken einzelner Töne, sondern alle Beziehungen und relativen Phasen zwischen ihnen, lässt sich die minimale Darstellung viel strenger festlegen. Man hat damit aber noch nicht erklärt, warum genau diese Partitur ausgewählt wurde. Das positive Quellfunktional ist Teil der Eingabe, solange es nicht aus einem tieferen Prinzip folgt.
\end{bild}

Die Rekonstruktion erzeugt nicht automatisch Lokalität, einen Zeitfluss, einen Diracoperator oder die Quantenmechanik aus $E_8$ allein. Komplexe Linearität und Positivität sind bereits starke Voraussetzungen. Bekannte operationelle Quantenrekonstruktionen benutzen ausdrücklich zusätzliche Informationsprinzipien, einschließlich einer Purifikationsforderung. \href{https://arxiv.org/abs/1011.6451}{Chiribella, D'Ariano und Perinotti [E7]}

\section{Mehrzeitdaten statt einer einzigen Dichtematrix}
Ein Prozesstensor beziehungsweise Quantenkamm ordnet einer Folge von Instrumenten gemeinsame Ergebniswahrscheinlichkeiten zu. Positivität und kausale Normierungsregeln sichern die Konsistenz. Eine tomographisch vollständige Eingriffsfamilie bestimmt einen endlichen Prozesstensor; sie bestimmt nicht automatisch eine eindeutige verborgene Umgebung oder sämtliche unendlichen Horizonte.

Eine Zeitstelle entfernt man durch Einsetzen einer Identitätsoperation. Das Summieren über Ergebnisse einer tatsächlich störenden Messung kann dagegen eine Dephasierung zurücklassen. Diese Unterscheidung ist für konsistente Mehrzeiterweiterungen wesentlich. \href{https://arxiv.org/abs/0904.4483}{Quantenkämme [E8]}, \href{https://arxiv.org/abs/1712.02589}{kausale Erweiterungssätze [E9]}

Die stärkste gemeinsame Arbeitsdefinition lautet damit
\begin{equation}
\mathcal U_{\mathrm{op}}=
\frac{\text{komponierbare Operationen + positive konsistente Prozessstatistik}}
{\text{Gleichheit unter allen erlaubten Experimenten}}.
\end{equation}
TFPT kann eine markierte Präsentation eines Teils davon sein. Um einen bestimmten physikalischen Prozess zu erhalten, müssen Komposition und Statistik aus derselben Quelle ausgewählt werden.

\section{Ein konkreter Auswahltest für offene Größen}
S05 schlägt eine nützliche endliche Vorgehensweise vor: Man fixiert zulässige Operationswörter bis Länge $L$, ihre Relationen, Symmetrien und die belegten Auslesungen. Die unbekannten Momente müssen eine positive Momentenmatrix bilden. Dann minimiert und maximiert man eine noch offene Ausgabe $\omega(O)$ über alle kompatiblen Modelle.

Bleibt ein Intervall, hat man eine genaue Nicht-Eindeutigkeit und eine messbare Frage. Kollabiert es, ist diese Ausgabe innerhalb der geprüften Relaxation festgelegt. Ein endlicher positiver Momentenkern ist noch keine Garantie einer konsistenten unendlichen Erweiterung; dafür sind weitere Bedingungen erforderlich. Der Vorteil ist dennoch klar: Eine Herkunftsvermutung wird durch konkrete Alternativen und Zeugen ersetzt.

Das Echo liefert bereits so einen Zeugen. Gleiche erste reduzierte Schritte erlauben verschiedene spätere Rückkehrwerte. Ein zukünftiger Ursprungskandidat muss die Registerausführung vorab festlegen oder die experimentelle Äquivalenz sämtlicher erlaubten Fortsetzungen beweisen.

\section{Die Zelle als Verknüpfungsamplitude}
Die in S03 vorgeschlagene Alternative verwendet
\begin{equation}
\Omega_{abcd}=\varepsilon_{abcd}/\sqrt{24}
\end{equation}
als lokalen Viererverknüpfungstensor. An orientierten Verbindungen werden Trägerräume mit ihren dualen Räumen kontrahiert. So verbindet man \emph{Amplituden}, ohne jede physische Zellreduktion als unverändert reines $\Omega$ festzuschreiben.

Die elementaren Identitäten lauten
\begin{align}
\sum_{bcd}\Omega_{abcd}\overline{\Omega_{a'bcd}}&=\frac14\delta_{aa'},\\
\sum_{cd}\Omega_{abcd}\overline{\Omega_{a'b'cd}}&=
\frac{\delta_{aa'}\delta_{bb'}-\delta_{ab'}\delta_{ba'}}{12}.
\end{align}
Sie reproduzieren die bekannten reduzierten Strukturen als Kompositionsregeln. Passende lokale Basiswechsel können sich an einer kontrahierten Verbindung aufheben; das motiviert eine Eichbeschreibung. Propagierende Eichbosonen folgen daraus noch nicht.

\begin{figure}[H]\centering
\begin{tikzpicture}
\node[box,circle,minimum size=11mm] (a) at (0,0) {$\varepsilon$};
\node[box,circle,minimum size=11mm] (b) at (3,0) {$\bar\varepsilon$};
\node[box,circle,minimum size=11mm] (c) at (3,2.2) {$\varepsilon$};
\node[box,circle,minimum size=11mm] (d) at (0,2.2) {$\bar\varepsilon$};
\draw[arr] (a)--(b);\draw[arr] (a)--(d);
\draw[arr] (c)--(b);\draw[arr] (c)--(d);
\draw[arr] (a)--(-1.2,0);\draw[arr] (a)--(0,-1.1);
\draw[arr] (4.2,0)--(b);\draw[arr] (3,-1.1)--(b);
\draw[arr] (c)--(4.2,2.2);\draw[arr] (c)--(3,3.3);
\draw[arr] (-1.2,2.2)--(d);\draw[arr] (0,3.3)--(d);
\node[openbox,text width=44mm] at (8,1.1) {Graph, Gewichte, Positivität und Registerregeln noch auszuwählen};
\end{tikzpicture}
\caption{Beispiel einer hypothetischen Tensorverklebung. Jeder Knoten besitzt vier Anschlüsse; $\varepsilon$ und sein dualer konjugierter Tensor wechseln sich ab. Innere Linien kontrahieren duale Anschlüsse, offene Linien bleiben als äußere Träger. Die Zeichnung zeigt keine schon abgeleitete Raumzeit.}
\end{figure}

Der eine $\mathrm{SU}(4)$-Tensor enthält noch nicht die vollständige markierte $E_8$-Struktur. $D_5$-Sektoren, $\Z_4$-Grade, Phasen und kompatible lokale Erweiterungen müssen in derselben Regel erscheinen. Der Vorschlag löst den Widerspruch unverändert reiner gekoppelter Zellen konzeptionell, aber nicht die Auswahl des Graphen oder Zustandsfunktionals.

\section{Teilsysteme aus Beziehungen erkennen}
\begin{bild}{Integrierter Arbeitsauftrag v1.6: die kleinste gemeinsame Ausführung}
Die nächsten Rekonstruktionsdaten sollten nicht weitere isolierte Zahlen sein. Gefordert ist eine einzige endliche Ausführung mit eindeutig benannten Eingangsmoden, nativen Operationen, behaltenen Hilfsregistern und vollständiger Rohstatistik. Ihre Übersetzung muss zugleich die markierte Matrixordnung, den signierten Vermittlungstensor und eine nichttriviale Ortsänderung erhalten. Die neuen Rechnungen liefern Teile dieser Übersetzung, aber noch kein solches Gesamtprotokoll. Ein Matrixwort mit derselben Zielwirkung zählt erst dann als nativ ausführbar, wenn auch seine linearen Kombinationen, Normierungen und Kontrollen aus erlaubten Ressourcen folgen.
\end{bild}


Paarweise kommutierende volle Matrixunteralgebren, die gemeinsam die Gesamtalgebra erzeugen, bestimmen im endlichen Fall eine Tensorproduktstruktur bis auf lokale Basiswechsel und Umbenennungen. Bei nichttrivialem Zentrum können stattdessen Sektoren übrigbleiben. Die Auswahl der privilegierten Unteralgebren aus den Wechselwirkungen ist der eigentliche inverse Schritt. \href{https://arxiv.org/abs/quant-ph/0308043}{Zanardi, Lidar und Lloyd [E10]}

Die Inversionsnotiz untersucht alle 105 Paarungen der acht vorgegebenen Qubits. Laut Quelle macht genau eine die sechs ursprünglichen Swaps zweikörperlich; vier Paarungen teilen das reine Rang-6-Marginalprofil. Der kombinierte Fingerabdruck identifiziert die Trägerpaarung in dieser endlichen Klasse. Das ist ein konkreter relationaler Test, keine Enumeration aller möglichen Faktorisierungen von $\C^{256}$.

\section{Was das Zustandsgate am Rotor/CAR-Ring prüft}
S08 untersucht einen neutralen Sektor aus 70 Materiemasken und Schleifenfluss, insgesamt 1190 Zuständen. In einer festgehaltenen lokalen Operatorliste wird über
\begin{equation}
\Gamma_{ij}=\operatorname{Re}\langle O_i\psi,O_j\psi\rangle
-\langle O_i\rangle\langle O_j\rangle
\end{equation}
der Generator rekonstruiert. Nach Entfernen trivialer Identitätsrichtungen reproduziert der Grundzustand die Koeffizienten $1/24$, $1/576$, $1/200$ und $383/96=M-\varepsilon_L$. Ein zusätzlicher angeregter Zustand verbessert die kleinste positive Kovarianzeigenzahl von etwa $1{,}3\cdot10^{-6}$ auf $1{,}7\cdot10^{-3}$.

Erweitert man die Liste um HH-Hop und $E^4$, bleibt der niedrige Zustand fast blind: Die relevanten Varianzen sind etwa $1{,}5\cdot10^{-13}$ und $2{,}8\cdot10^{-10}$. Er lebt überwiegend bei $|E|\le1$ und hat wenig H-Anteil. Zustände mit Energie in der Größenordnung $M$ oder unabhängige Quellregeln sind nötig. Das Verfahren $H\to\psi\to H$ prüft Identifizierbarkeit in einer gewählten Familie, nicht die unabhängige Herkunft dieser Familie. Die Originalringrechnung wurde hier nicht erneut ausgeführt.

\par\noindent\textbf{Neue Gegenprobe:} Kapitel~\ref{ch:shadow-quantum} konstruiert zwei Hamiltonoperatoren mit gleichen Energien und verschiedenen lokalen Grundzustandsantworten. Kapitel~\ref{ch:shadow-synthesis} ergänzt eine arithmetische Wegphase als konkrete gemischte Antwort.
